% !TeX program = lualatex
% Standalone research notebook for original catalogue problem 63.
% Compile twice: lualatex 63.tex
% Original entry retained; research subsection and [E] references added.
% No external bibliography, image, data, or style file is required.
\documentclass[11pt,a4paper]{article}
\usepackage[margin=24mm,headheight=14pt]{geometry}
\usepackage{fontspec}
\setmainfont{Latin Modern Roman}
\setsansfont{Latin Modern Sans}
\setmonofont{Latin Modern Mono}
\usepackage{amsmath,amssymb,mathtools,amsthm}
\usepackage{unicode-math}
\setmathfont{Latin Modern Math}
\usepackage{microtype}
\usepackage{enumitem,xurl,needspace}
\usepackage[dvipsnames]{xcolor}
\usepackage{hyperref}
\hypersetup{unicode=true,colorlinks=true,linkcolor=MidnightBlue,
 urlcolor=MidnightBlue,bookmarksnumbered=true,
 pdftitle={Research notebook - Problem 63},
 pdfauthor={Research attempt; not independently refereed}}
\usepackage{bookmark}
\usepackage{titlesec}
\titleformat{\section}{\Large\bfseries\raggedright}{\thesection}{0.7em}{}
\usepackage{fancyhdr}
\pagestyle{fancy}\fancyhf{}
\fancyhead[L]{\small\sffamily Research notebook}
\fancyhead[R]{\small\sffamily Problem 63}
\fancyfoot[L]{\scriptsize 27 September 2026}
\fancyfoot[C]{\thepage}
\fancyfoot[R]{\scriptsize Unrefereed research attempt}
\setlength{\parindent}{0pt}
\setlength{\parskip}{4pt plus 1pt}
\setlength{\emergencystretch}{3em}
\clubpenalty=10000
\widowpenalty=10000
\displaywidowpenalty=10000
\setcounter{secnumdepth}{1}
\setlist[itemize]{leftmargin=3em,itemsep=2pt,topsep=3pt,parsep=0pt,before=\small}
\allowdisplaybreaks
\newcommand{\N}{\mathbb{N}}
\newcommand{\Z}{\mathbb{Z}}
\newcommand{\Q}{\mathbb{Q}}
\newcommand{\R}{\mathbb{R}}
\newcommand{\C}{\mathbb{C}}
\newcommand{\F}{\mathbb{F}}
\newcommand{\A}{\mathbb{A}}
\newcommand{\PP}{\mathbb P}
\newcommand{\E}{\mathbb{E}}
\newcommand{\eps}{\varepsilon}
\newcommand{\dd}{\,\mathrm d}
\newcommand{\sref}[2]{\hyperlink{source:#1:#2}{[S#2]}}
\newcommand{\eref}[2]{\hyperlink{extra:#1:#2}{[E#2]}}
\newif\ifshowcatalogue
\showcataloguetrue
\newcommand{\cataloguescope}[1]{\ifshowcatalogue
 \paragraph{Exact scope in the supplied catalogue.}
 \begin{quote}\small #1\end{quote}\fi}
\newtheorem{notebooklemma}{Lemma}[section]
\newtheorem{notebookproposition}[notebooklemma]{Proposition}
\numberwithin{equation}{section}
\begin{document}
\setcounter{section}{62}
% ================================================================
% problemId: problem.nexp-versus-p-poly
% source releaseRank: 63; source releaseStatus: open
% exactTarget SHA-256: 72ebbcd8c1611e356f913428768ef5a22552850814192156bc6d0c6a51113926
\clearpage
\section{\texorpdfstring{NEXP versus P/poly}{NEXP versus P/poly}}\label{63}
\subsection{Definitions and mathematical statement}
Define $\mathsf{NEXP}=\bigcup_{k\ge1}\mathsf{NTIME}(2^{n^k})$, using nondeterministic machines with the indicated time bound on every branch. Membership in $\mathsf{P/poly}$ means that each input length has a polynomial-size Boolean circuit, with no uniform-generation condition. The question is
\[
 \mathsf{NEXP}\stackrel{?}{\subseteq}\mathsf{P/poly}.
\]
To refute containment, one needs one language in $\mathsf{NEXP}$ lacking every polynomial circuit-size bound, not merely a lower bound in a restricted circuit class.
\par\smallskip\textit{Formulation and concept references:} \sref{63}{1}.
\cataloguescope{Determine whether NEXP (nondeterministic exponential time) is contained in P/poly (the class of languages computable by polynomial-size Boolean circuit families).}
\subsection{Short English statement}
Could every nondeterministic exponential-time language nevertheless have polynomial-size nonuniform Boolean circuits?
% BEGIN RESEARCH INSERTION
\subsection{Research attempt}

\paragraph{Current frontier.}
Williams proves nondeterministic exponential-time lower bounds against polynomial-size $\mathsf{ACC}$, not unrestricted $\mathsf{P/poly}$ \eref{63}{1}, Theorem 1.1. Atserias--Buss--Müller establish a consistency result in bounded arithmetic; Atserias--Müller's April 2026 manuscript also concerns consistency and magnification, not an unconditional separation in the ordinary mathematical universe \eref{63}{2}, \eref{63}{3}. The qualifier ``consistent with a formal theory'' cannot be deleted. We audit direct diagonalization against the target's union-class quantifiers.

\paragraph{Attack route.}
Counting supplies many hard tables but not an efficient selector. Nondeterministically guessing a hard table still needs a certificate excluding all small circuits. The algorithm-to-lower-bound route instead needs a sufficiently fast general-circuit SAT algorithm. We push explicit selection until the circuit exponent's cost becomes visible, and test whether padding can repair it.

\paragraph{Attempt: fixed-exponent diagonalization with an actual algorithm.}
In a fixed finite fan-in-two basis, a topologically numbered circuit on $n$ inputs with at most $s$ gates has an encoding of length $O(s\log(n+s))$. The number of encodings has an explicit upper bound
\[
 M(n,s)\le2^{C_0s\log_2(n+s)}
\]
for an absolute $C_0$. Repetitions of the same function and invalid encodings cause no problem.
\begin{notebooklemma}\label{r63:fixed}
For every fixed positive integer $k$, there exists $L_k\in\mathsf{EXP}$ whose length-$n$ function has circuit size greater than $n^k$ for every sufficiently large $n$.
\end{notebooklemma}
\begin{proof}
Put $s=\lfloor n^k\rfloor$. For large $n$, $M(n,s)+1\le2^{2^n}$. Inspect the first $M(n,s)+1$ binary truth tables lexicographically. For each, enumerate all valid circuit encodings of size at most $s$ and evaluate them on all inputs. Select the first table differing from every circuit table. There are at most $M$ distinct circuit tables, so the search succeeds.

A direct implementation costs at most
\[
 (M(n,s)+1)M(n,s)2^n\operatorname{poly}(n+s)
        =2^{O(n^k\log n)}\le2^{O(n^{k+1})}.
\]
On input $x$, generate the canonical table for $n=|x|$ and return its $x$-entry. Handle the finitely many omitted lengths arbitrarily. This is a deterministic exponential-time language with the claimed bound.
\end{proof}
Only $M+1$ candidate tables, not all $2^{2^n}$ tables, are inspected. Candidate indices have $O(\log M)$ bits. The construction does not assume a nondeterministically checkable hardness certificate.

\paragraph{Where the union quantifier enters.}
The lemma establishes
\[
 \forall k\ \exists L_k\in\mathsf{EXP}:L_k\notin\mathsf{SIZE}(n^k),
\]
whereas the target needs one language defeating every fixed polynomial. The algorithm's time exponent depends on $k$. Replacing $k$ by an unbounded $k(n)$ gives no fixed $C$ certifying time $2^{n^C}$. Merely assigning different constructions to disjoint length intervals does not resolve this problem.

\paragraph{Attempted repair by padding.}
Suppose a core on $m$ variables is certified to require more than $m^k$ gates. Pad it to length $n$. Restricting padding variables in any padded circuit preserves the lower bound $m^k$, now measured against $n$. To obtain a fixed exponential-time budget from the preceding enumeration estimate, one requires
\[
 m^k\log m=O(n^C)
\]
for a fixed $C$. The same accounting then certifies only $m^k=O(n^C/\log m)$, rather than a superpolynomial lower bound in padded length. This is a barrier to this enumeration-and-padding calculation, not a theorem against every padding or magnification method.

\paragraph{A coherent selector would suffice.}
Suppose one deterministic machine prints a table $T_n$ in time $2^{n^C}$ for fixed $C$, and no circuit of size $n^{q(n)}$ computes it, where $q(n)\to\infty$. Its indexed-bit language is in $\mathsf{EXP}$ and not in $\mathsf{P/poly}$. Indeed a size $Bn^d$ family is eventually smaller than $n^{q(n)}$ because eventually $q(n)\ge d+1$ and $n>B$. This would give the stronger separation $\mathsf{EXP}\nsubseteq\mathsf{P/poly}$.

Nondeterministically guessing different hard tables on different branches does not define this language: accepting when some guessed table has a $1$ in a position gives their coordinatewise union, which need not be hard. A nondeterministic replacement needs an explicit coherence condition, such as identical output on all successful branches and at least one successful branch. No such selector is constructed.

\paragraph{A second route and its missing gain.}
Williams's general metatheorem uses a fixed sufficiently large inverse-polynomial improvement over exhaustive-search SAT time for every fixed polynomial circuit-size exponent to obtain nondeterministic exponential-time lower bounds \eref{63}{1}, Theorem 1.3. The direct evaluator costs $2^n\operatorname{poly}(n^k)$: its polynomial factor has the wrong sign. The present attempt supplies no improved general-circuit satisfiability algorithm.

\paragraph{Stress test.}
The selected table is the same for every evaluation at a given input length. Counting encodings rather than functions is a safe upper bound. Polynomial constants are absorbed only by passing to a larger fixed exponent, not by exchanging the quantifiers. $\mathsf{ACC}$ lower bounds do not cover arbitrary-depth circuits, and consistency statements do not construct the required selector. Finite arithmetic checks of the counting threshold are included in the script; they say nothing about a uniform unbounded-exponent selector.

\paragraph{Outcome.}
\textbf{Outcome: Exploratory approach with unresolved gap.}

A fully executable fixed-exponent diagonalizer and its padding tradeoff are supplied. They explain exactly why counting alone does not establish the union-class target. The unresolved input is a coherent fast selector or a sufficiently fast general-circuit SAT algorithm. The fixed-exponent result is a standard diagonalization phenomenon, reproduced and audited rather than presented as a new separation.

% END RESEARCH INSERTION
\subsection{Sources}
\begin{itemize}\raggedright
\item[\textnormal{[S1]}]\hypertarget{source:63:1}{} Complexity Zoo: NEXP entry.
\url{https://complexityzoo.net/Complexity_Zoo:N}.
{\small\textit{Catalogue formulation source.}}
\item[\textnormal{[S2]}]\hypertarget{source:63:2}{} On the Consistency of Circuit Lower Bounds for Non-Deterministic Time.
\url{https://arxiv.org/abs/2303.01016}.
{\small\textit{Catalogue research/status reference; not a proof endorsement.}}
\item[\textnormal{[S3]}]\hypertarget{source:63:3}{} From Godel incompleteness to the consistency of circuit lower bounds.
\url{https://arxiv.org/abs/2604.25251}.
{\small\textit{Catalogue research/status reference; not a proof endorsement.}}
\item[\textnormal{[S4]}]\hypertarget{source:63:4}{} Basic Circuit Complexity — Chi-Ning Chou.
\url{https://cnchou.github.io/notes/circuit.html}.
{\small\textit{Catalogue research/status reference; not a proof endorsement.}}

% BEGIN ADDED RESEARCH REFERENCES
\item[\textnormal{[E1]}]\hypertarget{extra:63:1}{} Ryan Williams, \emph{Non-Uniform ACC Circuit Lower Bounds}, Journal of the ACM 61(1), Article 2 (2014); author-hosted journal manuscript, Theorems 1.1 and 1.3.
\url{https://people.csail.mit.edu/rrw/acc-lbs-journal-final.pdf}.
{\small\textit{Restricted-circuit lower bound and the algorithmic metatheorem. Bibliographic journal data are taken from the author publication list, not its manuscript template footer. Consulted 27 September 2026.}}
\item[\textnormal{[E2]}]\hypertarget{extra:63:2}{} Albert Atserias, Sam Buss, and Moritz Müller, \emph{On the Consistency of Circuit Lower Bounds for Non-Deterministic Time}, arXiv:2303.01016v2 (26 August 2023).
\url{https://arxiv.org/abs/2303.01016v2}.
{\small\textit{A consistency theorem in bounded arithmetic, not an unconditional NEXP circuit lower bound. Consulted 27 September 2026.}}
\item[\textnormal{[E3]}]\hypertarget{extra:63:3}{} Albert Atserias and Moritz Müller, \emph{From Gödel incompleteness to the consistency of circuit lower bounds}, arXiv:2604.25251 (28 April 2026).
\url{https://arxiv.org/abs/2604.25251}.
{\small\textit{Consistency and proof-complexity magnification; recent manuscript, not a separation theorem. Consulted 27 September 2026.}}
% END ADDED RESEARCH REFERENCES
\end{itemize}

\end{document}
